\chapter*{Preface}
\addcontentsline{toc}{chapter}{Preface}

This book teaches logic through the act of getting Lean to accept or
reject a claim. It does not present logic as an abstract preliminary
to be endured before the real, practical work of programming begins;
the formal statement, the tactic script, and Lean's response
together constitute the lesson.

Each unit ends with a small number of experiments: typically one
successful proof and one or two countermodels demonstrating why a
plausible-looking alternative fails. A reader can learn as much from
discovering exactly why a converse refuses to typecheck as from
completing a valid proof. This pairing is not a pedagogical
decoration; it is the book's organizing method, carried through
consistently from the first chapter to the last.

The six chapters that follow form a complete, self-contained
introduction: propositions and connectives, natural deduction read
directly through Lean's tactics, the constructive/classical boundary
made observable rather than historical, quantifiers and their scope,
relations and orders, and elementary set theory treated as applied
logic. A further, more ambitious volume, \textit{Inference Under
Constraint}, extends this method into non-classical logic, modality,
and a broader admissibility-theoretic program; that volume assumes
this one but is not required to complete it. This book stands on its
own.

No prior experience with Lean, type theory, or formal proof is
assumed. Some comfort with ordinary mathematical notation is useful
but not essential — the notation is introduced as it is needed.
